\section{Introduction}
\label{sec:intro}

% black box testing is the ideal, but not the norm
%Early testing approaches focused validation effort on expected program behaviour~\cite{Howden1978FunctionalPT, richardson89}, rather than more commonly used approaches that focus on examining implementation structure~\cite{myers1979art}. 
Early testing approaches directed validation effort on expected program behaviour~\cite{Howden1978FunctionalPT, richardson89}, in contrast to the current more common approaches that examine implementation structure~\cite{myers1979art}. 
% Unfortunately, 
In practice, expected behaviours are often expressed through informal documentation, edge cases, and API contracts, making behavioural completeness difficult to evaluate systematically~\cite{hossain2025doc2oracllinvestigatingimpactdocumentation, 10.1109/ICST.2012.110}. 
As a result, software testing research and practice have largely converged on implementation-centric structural proxies for test quality~\cite{howden1976reliability, Budd1978Design, zhu1997software}, with
code coverage and mutation kill scores 
commonly being used to evaluate 
how complete a test suite is~\cite{papadakis2019mutation, 10.1145/2568225.2568278, 10.1145/2635868.2635929}.

% white box testing has downsides though
Unfortunately, structural test measures have introduced shortcomings that decouple test metrics from true behavioural intent. 
As these metrics validate correctness relative to the existing implementation, they are inherently blind to faults of omission: missing logic or unhandled requirements can still yield high structural scores~\cite{juristo2003functional, lawrance2005how}. 
%
These metrics also focus on code execution rather than behavioural validation; a test suite can maximize coverage or kill scores while lacking the robust oracles needed to validate expected behaviours, not just directly measuring the code as it was implemented~\cite{zhang2015assertions, petrovic2021does}. 
%
For instance, as we show in Section~\ref{sec:example}, a test achieving 100\% coverage can still leave an entire documented behaviour untested, leaving any regression against it undetected. 
%
Early software engineering theory warned that structural execution cannot guarantee that a program fulfills its expected behaviours~\cite{miller1963systematic, goodenough1975toward, hamlet1994foundations}; modern empirical evaluations confirm that actively optimizing for these code-centric thresholds yields highly optimized test suites with degraded fault-detection capabilities~\cite{10.1145/2568225.2568271, gay2015risks}.

% By contrast, developers encode a method's intended behaviours in its documentation, providing an abstract, implementation-independent specification of what the method should do~\cite{hossain2025doc2oracllinvestigatingimpactdocumentation}.  Though documentation is sometimes viewed as unreliable, inconsistencies with code are relatively infrequent ($\sim$10\%) and typically detectable~\cite{xu2025docprismlocalcategorizationexternal}, making it a reasonable source of truth for well-maintained libraries. The idea that structural metrics alone may be insufficient to capture test adequacy is not new. Fraser \textit{et. al.} argued that a test suite is adequate only if it is possible to infer an accurate model of a system's behaviour from its test execution~\cite{10.1109/ICST.2012.110}. Unfortunately, measuring whether a test suite validates documented behaviours has remained challenging because it requires linking natural language documentation to the behaviours evaluated by test cases~\cite{lee2025metamon}.

% need to soften:
The primary advantage of structural measures is that they are easy to collect as they map directly to the source code, which is amenable to automated analyses. 
This contrasts with expected behaviours, which are 
%To address these limitations, software testing theory has long motivated a shift toward evaluating adequacy against intended behaviours rather than structural targets. 
% In particular, Fraser \textit{et~al.} argue that a test suite is behaviourally adequate only if its execution traces are thorough enough to infer an accurate model of the system's intended behaviour~\cite{10.1109/ICST.2012.110}. 
%In practice, developers 
often encoded within natural language documentation, providing an abstract, implementation-independent description of what a method should do~\cite{hossain2025doc2oracllinvestigatingimpactdocumentation}, and are not currently directly measurable like coverage and kill scores. 
% While documentation is historically stereotyped as stale or unreliable, recent empirical evidence indicates that code-documentation inconsistencies are relatively infrequent ($\sim$10\%) and highly detectable~\cite{xu2025docprismlocalcategorizationexternal}, making it a viable source of truth for well-maintained libraries. 
While measuring whether a test suite validates documented behaviours has remained challenging because it requires linking natural language documentation to the behaviours evaluated by test cases~\cite{lee2025metamon}, 
recent advances in large language models (LLMs) now make this form of analysis feasible~\cite{Endres2023CanLL}. 
% i feel like these two sentences should be merged i meant the question and the core contribution
This paper therefore examines: \textit{To what extent can documentation-derived expected behaviours augment structural metrics in revealing  
inadequacies 
% behavioural gaps 
in existing test suites?}
We investigate this question empirically through four research questions around measuring behavioural coverage in software testing.
% The core contributions of this work empirically investigate four core research questions around measuring behavioural coverage in software testing.

% \textbf{RQ1: With what precision can expected behaviours be extracted from source code and documentation?} In this work, we therefore first investigate whether expected method-level behaviours can be extracted from source code and documentation with sufficiently high precision to support large-scale empirical analysis. To do this, we built \toolname{}, a tool that extracts expected behaviours from method documentation and code, maps behaviours to existing tests, and identifies which documented behaviours remain untested. Across ten popular open-source projects comprising 8,922 documented methods, \toolname{} surfaces 20,729 behaviours, 91.8\% of which are valid.

\textbf{RQ1 (Feasibility): With what precision can expected behaviours be extracted from source code and documentation?}
Before we can evaluate test adequacy, we must first determine if expected method-level behaviours can be extracted from natural language and code with sufficient precision for large-scale analysis. To this end, we introduce \toolname{}, a proof-of-concept that extracts expected behaviours, maps them to corresponding test cases, and identifies untested behaviours. Across ten popular open-source projects containing 8,922 documented methods, \toolname{} surfaces 20,729 distinct behaviours with a validated precision of 93.1\%.


%\textbf{RQ2: Do developer-written test suites exhibit behavioural gaps in practice?} To learn whether behavioural gaps exist within deployed systems, we further examine the test suites embedded in the projects examined in RQ1. We define a behavioural gap as a documented behaviour that is not validated by any existing test case. Across the ten projects, we find that 17.4\% of expected behaviours are not validated by the existing extensive test suites. Examining these gaps, we do not find that these missed behaviours were skipped because they are inherently untestable.

%\textbf{RQ2 (Prevalence): Do developer-written test suites exhibit behavioural gaps in practice?}
\textbf{RQ2 (Prevalence): To what extent do mature, developer-written test suites exhibit behavioural gaps?}
To understand the prevalence of behavioural gaps in real-world systems, we analyze the test suites of the same ten projects. We define a \textit{behavioural gap} as a documented behaviour that is not validated by any existing test case. Our analysis conservatively finds that at least 17.5\% of detected expected behaviours are not covered by these existing, well-maintained test suites. Furthermore, a qualitative inspection of these gaps indicates that these behaviours are rarely inherently untestable, suggesting a systematic shortfall in current validation practices.

% \textbf{RQ3: To what extent are automatically-generated test suites also prone to behavioural gaps?} To determine whether behavioural gaps are simply weaknesses introduced by human engineers we next looked at whether they also appear in automatically generated test suites. We examined whether two state-of-the-art automated test generation approaches, \textsc{EvoSuite}~\cite{10.1145/2025113.2025179} and \textsc{ASTER}~\cite{Pan2024ASTERNA} also exhibit behavioural gaps. Across four systems, we find that 26.8\% of behaviours are missed by the ASTER-generated suites and 20.4\% of behaviours are missed by the EvoSuite-generated suites, suggesting that the gaps missed in developer-written suites are also present in fully generated test suites.

%\textbf{RQ3 (Generalization): To what extent are automatically-generated test suites prone to behavioural gaps?}
\textbf{RQ3 (Generalization): Do behavioural gaps persist when test suites are automatically generated?}
To determine whether behavioural gaps are inherent to manual test authoring or a systemic limitation of structure-guided testing practices, we evaluate whether they persist in automatically generated test suites. We apply two state-of-the-art test generation approaches, \textsc{EvoSuite}~\cite{10.1145/2025113.2025179} and \textsc{ASTER}~\cite{Pan2024ASTERNA}, to four of our subject systems. We find that 27.1\% of behaviours are missed by \textsc{ASTER} and 20.6\% by \textsc{EvoSuite}. These results suggest that behavioural gaps are not just an artifact of human-written code, but a fundamental challenge for both human and automated testing strategies.

% \textbf{RQ4: Are behavioural gaps independent from coverage and kill score?} Finally, we hypothesize that behavioural gaps are not simply proxies for low line coverage or low mutation kill score and instead represent a complementary dimension of test adequacy that structural metrics cannot fully capture. Across TBD untested behaviours for developer-written tests, we find that 85.4\% are present in methods with $>90\%$ line coverage and 50.1\% are present in methods with $>90\%$ kill score, suggesting that untested behaviours can provide insight into test weaknesses, even for test suites that would be considered comprehensive by traditional empirical measures. % A Cohen's $h$ analysis shows behaviours have only a \textit{small} correlation with coverage, and a \textit{medium} correlation with kill score.

%\textbf{RQ4 (Implication): Are behavioural gaps independent of traditional structural metrics?}
\textbf{RQ4 (Implication): Do behavioural gaps alias traditional structural metrics?}
We hypothesize that behavioural gaps are not simply proxies for low coverage or mutation scores, but represent a distinct dimension of test adequacy. Across the untested behaviours identified in our developer-written suites, we find that the majority of untested behaviours occur in methods that already achieve high structural scores. Most fall in methods exceeding 90\% line coverage, and over half in methods exceeding 90\% mutation kill scores. %Further statistical analysis using Cohen's $h$ confirms that behavioural coverage exhibits only a small effect size with line coverage ($h\approx0.2$) and a medium effect size with kill scores ($h\approx0.4$). 
These findings indicate that behavioural gaps can identify validation weaknesses in code that is `comprehensively' tested, suggesting that behavioural gaps can complement analysis of structurally-strong suites.

% Across the untested behaviours identified in our developer-written suites, we find that 85.4\% occur within the TBD\% of methods already achieving $>90\%$ line coverage, and 50.1\% occur within the TBD\% of methods with $>90\%$ mutation kill scores. Further statistical analysis using Cohen's $h$ confirms that behavioural coverage exhibits only a small correlation with line coverage ($h<0.2$) and a medium correlation with kill scores ($h<0.5$). These findings indicate that behavioural gaps identify verification weaknesses in test suites that are otherwise considered ``comprehensive'' by traditional standards, suggesting that behavioural gaps can complement analysis of structurally-strong suites.


This paper makes the following key contributions:
\begin{itemize}
    \item \textbf{Empirical Analysis of Behavioural Gaps:} We conduct an empirical study across ten popular open-source Java libraries, demonstrating that significant behavioural gaps exist between the documented and validated behaviours of real-world methods, even in suites with high structural coverage.
    
    \item \textbf{Generalization to Automated Testing:} We present a direct comparison against two state-of-the-art automated test generation approaches, confirming that these behavioural gaps persist even in fully-generated test suites.

    \item \textbf{Behavioural Coverage as a Distinct Dimension:} We show that behavioural gaps do not alias line coverage or mutation kill score. Gaps persist in 38.2\% of methods with perfect line coverage and 29.6\% with a perfect mutation kill score, establishing behavioural coverage as a complementary dimension of test adequacy. %rather than a restatement of existing metrics.
    
    \item \textbf{Tool \& Artifacts:} We provide a publicly available replication package, including \toolname{} and all identified behavioural gaps, to enable researchers to reproduce our results and extend this analysis to other software systems.\footnote{\toolname{}: \url{https://anonymous.4open.science/r/behavioural-coverage}}
\end{itemize}

\noindent Section~\ref{sec:example} provides a motivating example showing how behaviours are missed, despite being covered. Section~\ref{sec:related} describes related work. \toolname{} is introduced in Section~\ref{sec:approach}. Section~\ref{sec:eval} describes empirical study and its findings.
% RTH: decide if discussion needs to come back
Section~\ref{sec:summary} concludes.
